\begin{table}[htbp]
\centering
\caption{Standardized Effect Sizes}
\label{tab:sde}
\begin{tabular}{lcccccc}
\toprule
Outcome & $\hat{\beta}$ & SE & SD($Y$) & SDE & SE(SDE) & Classification \\
\midrule
\multicolumn{7}{l}{\textit{Panel A: Pooled}} \\[3pt]
$\ln(\text{Patents})$ & 0.2043 & 0.1172 & 0.8456 & 0.2415 & 0.1386 & Large positive \\
Claims/patent & -0.0242 & 0.3567 & 2.2436 & -0.0108 & 0.1590 & Small negative \\
Forward citations & 0.9665 & 0.9198 & 4.7450 & 0.2037 & 0.1938 & Large positive \\
\midrule
\multicolumn{7}{l}{\textit{Panel B: Heterogeneous (sample splits)}} \\[3pt]
$\ln(\text{Pat.})$, Semicond. & 0.2766 & 0.0892 & 0.8456 & 0.3271 & 0.1055 & Large positive \\
$\ln(\text{Pat.})$, Optoelec. & 0.2685 & 0.1515 & 0.8456 & 0.3175 & 0.1792 & Large positive \\
\bottomrule
\end{tabular}
\begin{tablenotes}
\small
\item \textit{Notes:} \textbf{Country:} Taiwan. \textbf{Research question:} Did the 2010 transition from sector-targeted R\&D tax credits (up to 35\% for strategic industries under the SUI) to a uniform 15\% credit (IIA) reduce patenting effort in previously favored technology classes at the USPTO? \textbf{Policy mechanism:} The IIA replaced the SUI's sector-specific enhanced credits with a uniform 15\% credit for all industries, imposing an effective 20 percentage-point credit reduction on designated strategic sectors including semiconductors, optoelectronics, communications, and IT hardware. \textbf{Outcome definition:} Annual count of utility patent applications filed at the USPTO by Taiwan-based assignees, aggregated at the USPC technology class level; claims per patent measures patent scope; forward citations count subsequent US patents citing each patent. \textbf{Treatment:} Binary; 22 USPC classes corresponding to SUI-designated strategic industries are treated after January 2010. \textbf{Data:} PatentsView bulk download, 2003--2013, USPC class $\times$ year level, approximately 190,000 utility patents across 50 technology classes and 11 years for three countries (Taiwan, Israel, South Korea). \textbf{Method:} Two-way fixed effects difference-in-differences with USPC class and year fixed effects; standard errors clustered at the USPC class level; heterogeneity via sample splits between semiconductor and optoelectronics sub-sectors. \textbf{Sample:} Top 50 USPC classes by Taiwan patent volume; utility patents only. SDE $= \hat{\beta} / \text{SD}(Y)$ where SD($Y$) is the pre-treatment standard deviation. Classification refers to magnitude, not statistical significance: Large ($|$SDE$| > 0.15$), Moderate ($0.05$--$0.15$), Small ($0.005$--$0.05$), Null ($< 0.005$).
\end{tablenotes}
\end{table}
